% GENERATED FILE. Edit canonical lessons, then run npm run generate:books.
% canonical-source-hash: fnv1a64:f3dd2af9951105f8
% canonical-lessons: GE-C60-gesund, GE-C60-durst, GE-C60-hunger, GE-C60-traurig, GE-C60-gluecklich

\chapter{Healthy, Thirst, Hunger, Sad, Happy}
\label{ch:ge-healthy-thirst-hunger-sad-happy}
\hlchaptermodality{60}

\begin{chapteropening}
\textbf{By the end of this chapter:} \emph{I can say healthy, thirst, hunger, sad and happy.}

Complete the last lesson of chapter 60: I can say healthy, thirst, hunger, sad and happy.
\end{chapteropening}

\section[gesund]{gesund — healthy}

\label{lesson:GE-C60-gesund}



\begin{quote}
Before the new one: say the German for a hospital, then the German for ill.
\end{quote}

\subsection*{You'll want to know: gesund}
\textbf{gesund} — \textquotedblleft{}healthy\textquotedblright{}.

\begin{cousinweb}
\textbf{gesund} $\leftarrow$ Old High German \emph{gisunt} $\leftarrow$ Proto-Germanic \emph{*gasundaz}, \textquotedblleft{}sound, whole, healthy\textquotedblright{} — inherited, a cousin of English \emph{sound}, as in \textquotedblleft{}safe and sound\textquotedblright{}. (EtymDB 2.1)
\end{cousinweb}

\begin{grammarlens}[title={What you've built}]
One more word for this chapter.
\end{grammarlens}

\subsection*{Your turn}
\begin{itemize}
\raggedright
  \item \emph{Say it:} \emph{gesund}
  \item \emph{Say it:} \emph{gesund}, once more
  \item \emph{Say it:} say \emph{gesund}
  \item \emph{Recall:} say the German for a hospital, then the German for ill, then say \emph{gesund} again
\end{itemize}

\begin{tcolorbox}[breakable,colback=teal!4,colframe=teal!35!black,title={Before you move on}]
What does gesund mean? (\textquotedblleft{}Healthy\textquotedblright{}.) Say it once more.
\end{tcolorbox}
\section[der Durst]{der Durst — thirst}

\label{lesson:GE-C60-durst}



\begin{quote}
Before the new one: say the German for ill, then the German for healthy.
\end{quote}

\subsection*{You'll want to know: der Durst}
\textbf{der Durst} — \textquotedblleft{}thirst\textquotedblright{}.

\begin{cousinweb}
\textbf{Durst} $\leftarrow$ Old High German \emph{durst} $\leftarrow$ Proto-Germanic \emph{*þurstuz}, \textquotedblleft{}dryness, thirst\textquotedblright{} — inherited, the same word as English \emph{thirst}. (EtymDB 2.1)
\end{cousinweb}

\begin{grammarlens}[title={What you've built}]
One more word for this chapter.
\end{grammarlens}

\subsection*{Your turn}
\begin{itemize}
\raggedright
  \item \emph{Say it:} \emph{der Durst}
  \item \emph{Say it:} \emph{der Durst}, once more
  \item \emph{Say it:} say \emph{der Durst}
  \item \emph{Recall:} say the German for ill, then the German for healthy, then say \emph{der Durst} again
\end{itemize}

\begin{tcolorbox}[breakable,colback=teal!4,colframe=teal!35!black,title={Before you move on}]
What does der Durst mean? (\textquotedblleft{}Thirst\textquotedblright{}.) Say it once more.
\end{tcolorbox}
\section[der Hunger]{der Hunger — hunger}

\label{lesson:GE-C60-hunger}



\begin{quote}
Before the new one: say the German for healthy, then the German for thirst.
\end{quote}

\subsection*{You'll want to know: der Hunger}
\textbf{der Hunger} — \textquotedblleft{}hunger\textquotedblright{}.

\begin{cousinweb}
\textbf{Hunger} $\leftarrow$ Old High German \emph{hungar} $\leftarrow$ Proto-Germanic \emph{*hungruz} — inherited, the same word as English \emph{hunger}. (EtymDB 2.1)
\end{cousinweb}

\begin{grammarlens}[title={What you've built}]
One more word for this chapter.
\end{grammarlens}

\subsection*{Your turn}
\begin{itemize}
\raggedright
  \item \emph{Say it:} \emph{der Hunger}
  \item \emph{Say it:} \emph{der Hunger}, once more
  \item \emph{Say it:} say \emph{der Hunger}
  \item \emph{Recall:} say the German for healthy, then the German for thirst, then say \emph{der Hunger} again
\end{itemize}

\begin{tcolorbox}[breakable,colback=teal!4,colframe=teal!35!black,title={Before you move on}]
What does der Hunger mean? (\textquotedblleft{}Hunger\textquotedblright{}.) Say it once more.
\end{tcolorbox}
\section[traurig]{traurig — sad}

\label{lesson:GE-C60-traurig}



\begin{quote}
Before the new one: say the German for thirst, then the German for hunger.
\end{quote}

\subsection*{You'll want to know: traurig}
\textbf{traurig} — \textquotedblleft{}sad\textquotedblright{}.

\begin{cousinweb}
\textbf{traurig} $\leftarrow$ Middle High German \emph{trūrec} $\leftarrow$ Old High German \emph{trūrag} (EtymDB 2.1) — inherited. A link to English \emph{dreary} is often suggested, but it is uncertain.
\end{cousinweb}

\begin{grammarlens}[title={What you've built}]
One more word for this chapter.
\end{grammarlens}

\subsection*{Your turn}
\begin{itemize}
\raggedright
  \item \emph{Say it:} \emph{traurig}
  \item \emph{Say it:} \emph{traurig}, once more
  \item \emph{Say it:} say \emph{traurig}
  \item \emph{Recall:} say the German for thirst, then the German for hunger, then say \emph{traurig} again
\end{itemize}

\begin{tcolorbox}[breakable,colback=teal!4,colframe=teal!35!black,title={Before you move on}]
What does traurig mean? (\textquotedblleft{}Sad\textquotedblright{}.) Say it once more.
\end{tcolorbox}
\section[glücklich]{glücklich — happy}

\label{lesson:GE-C60-gluecklich}



\begin{quote}
Before the new one: say the German for hunger, then the German for sad.
\end{quote}

\subsection*{You'll want to know: glücklich}
\textbf{glücklich} — \textquotedblleft{}happy\textquotedblright{}.

\begin{cousinweb}
\textbf{glücklich} is the word for \textquotedblleft{}luck\textquotedblright{}, Middle High German \emph{gelücke} (EtymDB 2.1), + \emph{-lich}. English \emph{luck} is the same word, taken from Low German or Dutch; where \emph{gelücke} itself came from is uncertain.
\end{cousinweb}

\begin{grammarlens}[title={What you've built}]
One more word for this chapter.
\end{grammarlens}

\subsection*{Your turn}
\begin{itemize}
\raggedright
  \item \emph{Say it:} \emph{glücklich}
  \item \emph{Say it:} \emph{glücklich}, once more
  \item \emph{Say it:} say \emph{glücklich}
  \item \emph{Recall:} say the German for hunger, then the German for sad, then say \emph{glücklich} again
\end{itemize}

\begin{tcolorbox}[breakable,colback=teal!4,colframe=teal!35!black,title={Before you move on}]
What does glücklich mean? (\textquotedblleft{}Happy\textquotedblright{}.) Say it once more.
\end{tcolorbox}
